\documentclass[12pt, a4paper]{article}

\usepackage[top=2.5cm, bottom=2.5cm, left=3cm, right=3cm]{geometry}
\usepackage{mdframed}
\usepackage[utf8]{inputenc}
\usepackage[T1]{fontenc}
\usepackage[spanish]{babel}
\usepackage{amsmath}
\usepackage{float}
\usepackage[margin=10pt, font=small, labelfont=bf]{caption}
\usepackage{booktabs}
\usepackage{graphicx}
\usepackage{xcolor}
\usepackage{hyperref}
\usepackage{natbib}
\usepackage{longtable}
\usepackage{setspace}
\usepackage{titlesec}
\usepackage{fancyhdr}

\setlength{\textfloatsep}{18pt plus 2pt minus 4pt}
\setlength{\floatsep}{12pt plus 2pt minus 2pt}
\setlength{\intextsep}{14pt plus 2pt minus 2pt}
\setstretch{1.15}
\setlength{\parindent}{0pt}
\setlength{\parskip}{8pt}

\titleformat{\section}
  {\large\bfseries}{\thesection.}{0.5em}{}
  [\vspace{2pt}\hrule\vspace{4pt}]
\titleformat{\subsection}
  {\normalsize\bfseries}{\thesubsection.}{0.5em}{}
\titlespacing{\section}{0pt}{18pt}{8pt}
\titlespacing{\subsection}{0pt}{14pt}{6pt}

\hypersetup{colorlinks=true, linkcolor=blue!60!black}

\pagestyle{fancy}
\fancyhf{}
\rhead{\small }
\lhead{\small \leftmark}
\cfoot{\thepage}
\renewcommand{\headrulewidth}{0.4pt}

\geometry{margin=1in}
\graphicspath{{Paper/figures/}{figures/}}
\newcommand{\tablepath}{Paper/tables/}
\newcommand{\figurelang}{es}
\newif\ifpolicyreport
\policyreporttrue

\title{Mercado laboral, ingresos y brechas en Colombia, 2010--2025\\
Un análisis descriptivo a partir de la GEIH}
\author{Centro Javeriano de Competitividad}
\date{\today}

\begin{document}

\maketitle
\begin{mdframed}[
  linewidth=1pt,
  linecolor=gray!60,
  backgroundcolor=gray!6,
  innertopmargin=12pt,
  innerbottommargin=12pt,
  innerleftmargin=14pt,
  innerrightmargin=14pt
]
\section*{Resumen ejecutivo}
\textbf{Este informe caracteriza la evolución del mercado laboral colombiano entre 2010 y 2025 a partir de los microdatos de la Gran Encuesta Integrada de Hogares (GEIH).} El documento examina los cambios en ocupación, ingresos laborales por hora y horas trabajadas, y luego desagrega esas dinámicas por sexo, formalidad, nivel educativo, posición ocupacional, departamento, sector económico y tamaño de empresa.

El mercado laboral colombiano terminó el período con más ocupados y mayores ingresos reales por hora. Sin embargo, los avances agregados conviven con brechas persistentes: la informalidad continúa siendo elevada, los retornos a la educación superior siguen marcando diferencias importantes, las desigualdades territoriales son profundas y la composición sectorial del empleo condiciona fuertemente las trayectorias de ingreso.

La motivación del informe es ofrecer una lectura detallada de quiénes han mejorado sus condiciones laborales, quiénes permanecen rezagados y qué dimensiones deberían orientar una agenda de política laboral, formación de capital humano y desarrollo productivo.
\end{mdframed}

\section{Mercado laboral, ingresos y brechas en Colombia}

\begin{figure}[H]
  \centering
  \includegraphics[width=0.9\textwidth]{Paper/figures/fig8Informe.png}
  \caption{Evolución general de los ocupados}
  \label{fig:1Informe}
\end{figure}


\textbf{El mercado laboral colombiano ha experimentado una expansión sostenida en el número de ocupados, alcanzando en 2025 su nivel histórico más alto con 23,0 millones de trabajadores.} Entre 2010 y 2025, la ocupación creció a una tasa anualizada del 2,09\%, lo que representa la incorporación de aproximadamente 6,2 millones de nuevos trabajadores en quince años. Este crecimiento, sin embargo, no fue lineal: el período 2014--2018 mostró una relativa estabilización en torno a los 18,5 millones de ocupados, seguida de un máximo prepandemia de 19,6 millones en 2019. La crisis sanitaria de 2020--2021 interrumpió abruptamente esta tendencia, con una caída hasta 17,6 millones de ocupados en 2021, el nivel más bajo de la década. La recuperación posterior fue notable: entre 2021 y 2025 el mercado laboral no solo recuperó el terreno perdido, sino que lo superó ampliamente, añadiendo 5,4 millones de empleos en apenas cuatro años.

\begin{figure}[H]
  \centering
  \includegraphics[width=0.9\textwidth]{Paper/figures/fig1Informe.png}
  \caption{Evolución general del ingreso laboral}
  \label{fig:2Informe}
\end{figure}

\textbf{Los ingresos laborales por hora han crecido en términos reales a lo largo del período analizado, aunque con una aceleración particularmente marcada en los últimos años.} Medido en pesos constantes de 2025, el ingreso laboral promedio por hora pasó de \$8.039 en 2010 a \$10.453 en 2025, lo que equivale a un crecimiento anualizado del 1,77\% real anual. Durante el período 2014--2019, los ingresos mostraron un comportamiento estancado, oscilando entre \$8.320 y \$8.839, sin una tendencia clara de mejora. A partir de 2021, coincidiendo con la recuperación postpandemia y los ajustes al salario mínimo, se observa una inflexión positiva sostenida que elevó el ingreso por hora en más de \$1.700 pesos reales en cuatro años, el avance más significativo de todo el período.

\textbf{El mercado laboral colombiano muestra en 2025 una configuración simultáneamente más grande y mejor remunerada que en cualquier momento de los últimos quince años, lo que sugiere una mejora estructural en las condiciones de empleo.} La expansión conjunta del número de ocupados y del ingreso real por hora indica que el crecimiento reciente no se ha producido a costa de la calidad del empleo, al menos en términos agregados. No obstante, esta lectura positiva debe matizarse: el crecimiento anualizado del ingreso por hora (1,77\%) es modesto en perspectiva histórica, y la magnitud del shock de 2020--2021 evidencia la alta vulnerabilidad del mercado laboral colombiano ante choques externos. El análisis desagregado por categorías de formalidad, sector y perfil del trabajador, que se desarrolla en las secciones siguientes, permitirá establecer si estos avances agregados se traducen en mejoras homogéneas o si encubren dinámicas heterogéneas entre distintos grupos de la población ocupada.

\subsection{Análisis por sexo}

\begin{figure}[H]
  \centering
  \includegraphics[width=0.9\textwidth]{Paper/figures/fig3Informe.png}
  \caption{Evolución de la composición laboral por sexo}
  \label{fig:3Informe}
\end{figure}

\textbf{La ocupación en Colombia sigue siendo predominantemente masculina, aunque la participación de las mujeres en el mercado laboral ha registrado una ganancia gradual de 1,3 puntos porcentuales entre 2010 y 2025.} En 2010, los hombres representaban el 60,1\% de los ocupados frente al 39,9\% de las mujeres; para 2025 esta proporción se sitúa en 58,8\% y 41,2\%, respectivamente. La tendencia es estable a lo largo del período, sin inflexiones abruptas, lo que sugiere que el avance en la participación femenina responde a un proceso estructural de largo plazo antes que a choques coyunturales. No obstante, la brecha de 17,6 puntos porcentuales que persiste en 2025 indica que la convergencia entre sexos en el acceso al empleo, aunque real, avanza a un ritmo lento.

\begin{figure}[H]
  \centering
  \includegraphics[width=0.9\textwidth]{Paper/figures/fig2Informe.png}
  \caption{Evolución del ingreso laboral por hora según sexo}
  \label{fig:4Informe}
\end{figure}

\textbf{Las mujeres han logrado cerrar parcialmente la brecha salarial por hora frente a los hombres, e incluso la han revertido en 2025, percibiendo un ingreso por hora levemente superior al masculino.} En pesos constantes de 2025, el ingreso por hora de los hombres pasó de \$8.176 en 2010 a \$10.400 en 2025, con un crecimiento anualizado del 1,62\%; el de las mujeres, en cambio, partió de un nivel inferior de \$7.832 y alcanzó \$10.527, con un crecimiento anualizado del 1,99\%. Esta dinámica diferencial implicó que las trayectorias de ambos sexos, que durante la mayor parte del período 2010--2019 mostraban al ingreso masculino por encima, convergieran progresivamente hasta cruzarse en los últimos años. Este resultado es llamativo y merece interpretarse con cautela: puede reflejar tanto mejoras reales en la calidad del empleo femenino como cambios en la composición sectorial u ocupacional de lastrabajadoras, aspectos que se examinarán en secciones posteriores.

\begin{figure}[H]
  \centering
  \includegraphics[width=0.9\textwidth]{Paper/figures/fig4Informe.png}
  \caption{Evolución de horas semanales trabajadas según sexo}
  \label{fig:5Informe}
\end{figure}

\textbf{Los hombres trabajan consistentemente más horas semanales que las mujeres, aunque ambos sexos muestran una tendencia decreciente en la jornada laboral a lo largo del período.} En 2010, los hombres registraban una jornada promedio de 50,2 horas semanales, que se redujo a 46,0 horas en 2025, con un decrecimiento anualizado del 0,58\%; las mujeres, por su parte, pasaron de 41,8 a 40,5 horas, con un decrecimiento más moderado del 0,2\% anual. La reducción de la jornada masculina es estadísticamente más pronunciada y sostenida, lo que ha contribuido a acercar las intensidades laborales entre sexos, aunque en 2025 persiste una diferencia de 5,5 horas semanales a favor de los hombres. Este patrón es relevante para interpretar los ingresos totales: aun cuando el ingreso por hora femenino converge con el masculino, la mayor jornada de los hombres implica que las brechas en ingresos laborales mensuales o anuales siguen siendo más amplias de lo que la comparación horaria sugiere.

\subsection{Análisis por formalidad}

\begin{figure}[H]
  \centering
  \includegraphics[width=0.9\textwidth]{Paper/figures/fig7Informe.png}
  \caption{Composición del empleo por formalidad}
  \label{fig:6Informe}
\end{figure}

\textbf{La informalidad laboral continúa siendo el rasgo estructural más pronunciado del mercado laboral colombiano, aunque entre 2010 y 2025 se registra una reducción significativa de 12,5 puntos porcentuales en la participación de los trabajadores informales.} En 2010, el 67,6\% de los ocupados era informal y apenas el 32,4\% cotizaba al sistema de pensiones; para 2025, estas proporciones se sitúan en 55,1\% y 44,9\%, respectivamente. La tendencia hacia la formalización ha sido sostenida y prácticamente ininterrumpida a lo largo del período, sin reversiones importantes ni siquiera durante el choque de 2020--2021, lo que sugiere que la recuperación postpandemia incorporó proporcionalmente más empleo formal que informal. A pesar de este avance, más de la mitad de los trabajadores colombianos permanece en la informalidad en 2025, lo que limita su acceso a la protección social y a los beneficios asociados al empleo formal.

\begin{figure}[H]
  \centering
  \includegraphics[width=0.9\textwidth]{Paper/figures/fig5Informe.png}
  \caption{Ingreso laboral por hora según formalidad}
  \label{fig:7Informe}
\end{figure}

\textbf{La brecha salarial entre trabajadores formales e informales es amplia y persistente: en 2025, un trabajador formal percibe por hora más del doble que uno informal, en términos reales.} El ingreso por hora de los formales pasó de \$12.824 en 2010 a \$15.241 en 2025, con un crecimiento anualizado del 1,16\%; el de los informales, en cambio, partió de \$5.575 y alcanzó \$6.271, creciendo apenas al 0,79\% anual. Las trayectorias de ambos grupos son marcadamente divergentes: mientras el ingreso formal muestra una tendencia creciente con aceleración visible a partir de 2021, el ingreso informal permanece prácticamente estancado durante todo el período 2010--2024, con una leve mejora solo en el último año. Este patrón indica que los beneficios del crecimiento económico y de los ajustes salariales recientes se han concentrado desproporcionadamente en los trabajadores formales.

\begin{figure}[H]
  \centering
  \includegraphics[width=0.9\textwidth]{Paper/figures/fig9Informe.png}
  \caption{Tasa de informalidad laboral}
  \label{fig:8Informe}
\end{figure}

\textbf{La tasa de informalidad ha descendido de manera continua durante quince años, acumulando una caída de 12,5 puntos porcentuales, con una aceleración particularmente marcada en el período 2021--2025.} Definida como la proporción de ocupados que no cotiza a pensión, la informalidad pasó del 67,6\% en 2010 al 55,1\% en 2025. La reducción fue gradual entre 2010 y 2019, con una caída promedio de aproximadamente 0,7 puntos porcentuales por año; sin embargo, entre 2021 y 2025 el ritmo se aceleró notoriamente, con descensos de hasta 1,4 puntos porcentuales anuales. Pese a este progreso, Colombia mantiene aún una tasa de informalidad superior al 55\%, que la ubica entre los mercados laborales con mayor prevalencia de empleo no protegido en América Latina, lo que representa un desafío estructural de primer orden para la política pública.

\begin{figure}[H]
  \centering
  \includegraphics[width=0.9\textwidth]{Paper/figures/fig6Informe.png}
  \caption{Razón del ingreso formal sobre el informal}
  \label{fig:9Informe}
\end{figure}

\textbf{La razón entre el ingreso por hora de los trabajadores formales e informales se ha ampliado entre 2010 y 2025, lo que evidencia que la formalización del empleo no ha ido acompañada de una convergencia salarial entre ambos segmentos.} En 2010, un trabajador formal ganaba en promedio 2,30 veces lo que ganaba uno informal por hora; en 2025 esa razón asciende a 2,43 veces. La evolución no es lineal: entre 2010 y 2019 la brecha se mantuvo relativamente estable en torno a 2,2 veces, pero a partir de 2020 se amplió de forma sostenida hasta alcanzar un máximo de 2,63 veces en 2023, para luego moderarse ligeramente. Este comportamiento sugiere que, si bien más trabajadores acceden al empleo formal, la prima salarial asociada a la formalidad se ha incrementado, lo que hace más costosa en términos relativos la permanencia en la informalidad y refuerza la importancia de acelerar la transición hacia el empleo protegido.

\subsection{Análisis por nivel educativo}

\begin{figure}[H]
  \centering
  \includegraphics[width=0.9\textwidth]{Paper/figures/fig12Informe.png}
  \caption{Composición del empleo por nivel educativo}
  \label{fig:10Informe}
\end{figure}

\textbf{La fuerza laboral colombiana ha experimentado una transformación educativa profunda entre 2010 y 2025, con un desplazamiento sostenido desde los niveles de menor cualificación hacia la educación superior.} La participación de trabajadores con educación superior o universitaria creció 11 puntos porcentuales, pasando del 22,8\% al 33,7\% del total de ocupados, consolidándose como el grupo con mayor crecimiento relativo del período. En sentido contrario, los trabajadores con básica primaria redujeron su participación en 12,6 puntos porcentuales, del 30,7\% al 18,1\%, y los de ningún nivel o preescolar cayeron 2,4 puntos hasta el 2,9\%. La educación básica secundaria o media, por su parte, ganó 4 puntos porcentuales hasta representar el 45,3\% de los ocupados, consolidándose como el nivel educativo más frecuente en el mercado laboral. Este proceso de upgrading educativo de la fuerza laboral constituye uno de los cambios estructurales más relevantes del período y tiene implicaciones directas sobre la productividad agregada y la composición de los ingresos laborales.

\begin{figure}[H]
  \centering
  \includegraphics[width=0.9\textwidth]{Paper/figures/fig10Informe.png}
  \caption{Ingreso laboral por hora según nivel educativo}
  \label{fig:11Informe}
\end{figure}

\textbf{El retorno educativo en el mercado laboral colombiano es pronunciado y se mantiene en 2025: un trabajador con educación superior percibe por hora más del doble que uno con básica secundaria, y más de tres veces lo que gana uno sin ningún nivel educativo.} En pesos constantes de 2025, el ingreso por hora de los trabajadores con educación superior asciende a \$18.013, frente a \$7.105 para quienes tienen básica secundaria o media, \$5.709 para quienes tienen básica primaria, y \$4.338 para quienes no tienen ningún nivel educativo. Comparado con 2010, todos los grupos registran mejoras en términos absolutos, pero la distancia entre el nivel superior y los demás se mantiene amplia. La brecha entre el grupo de mayor y menor educación pasó de \$13.466 en 2010 a \$13.675 en 2025, lo que indica que, en valores absolutos, la desigualdad salarial por educación no se ha reducido durante el período.

\begin{figure}[H]
  \centering
  \includegraphics[width=0.9\textwidth]{Paper/figures/fig11Informe.png}
  \caption{Ingreso laboral por hora según nivel educativo}
  \label{fig:12Informe}
\end{figure}

\textbf{En términos de crecimiento relativo, los trabajadores con menor nivel educativo son quienes han visto aumentar más rápidamente su ingreso por hora real entre 2010 y 2025, mientras que los universitarios registran el avance más modesto.} Con base 2010 igual a 100, los trabajadores sin ningún nivel educativo alcanzaron un índice de 121,4 en 2025, equivalente a un crecimiento anualizado del 1,3\%; los de básica primaria llegaron a 120,7, con un 1,26\% anual; y los de básica secundaria o media a 116,5, con un 1,02\% anual. En contraste, los trabajadores con educación superior registraron el crecimiento más bajo, con un índice de apenas 105,7 y una tasa anualizada del 0,37\%. Este patrón de convergencia relativa es alentador desde la perspectiva de la equidad, pues sugiere que los salarios de los trabajadores menos cualificados han crecido proporcionalmente más; sin embargo, dado que parten de niveles absolutos muy inferiores, la reducción de la brecha salarial en términos reales sigue siendo marginal y el diferencial por educación permanece como uno de los determinantes más potentes de la desigualdad de ingresos en Colombia.

\subsection{Análisis por posición ocupacional}

\begin{figure}[H]
  \centering
  \includegraphics[width=0.9\textwidth]{Paper/figures/fig15Informe.png}
  \caption{Composición del empleo por posición ocupacional}
  \label{fig:13Informe}
\end{figure}

\textbf{La estructura ocupacional del mercado laboral colombiano se ha desplazado hacia el empleo asalariado privado, que gana 9 puntos porcentuales entre 2010 y 2025, en detrimento del trabajo por cuenta propia y de las posiciones patronales.} En 2010, las dos categorías dominantes eran cuenta propia (45,1\%) y empleado particular (36,7\%), configurando un mercado laboral con una base amplia de autoempleo; para 2025, esta relación se invierte y el empleo particular asciende al 45,7\% mientras que el cuenta propia se reduce al 41,6\%. Las demás posiciones registran contracciones moderadas: patrón o empleador cae 2,7 puntos hasta el 2,3\%, empleado de gobierno retrocede 0,8 puntos hasta el 3,7\%, y servicio doméstico pierde 0,8 puntos hasta situarse en el 3,1\%. Este reordenamiento refleja un proceso de asalarización progresiva de la fuerza laboral que, leído junto con la reducción de la informalidad documentada en la sección anterior, sugiere una transformación estructural hacia formas de empleo más protegidas y estables.

\begin{figure}[H]
  \centering
  \includegraphics[width=0.9\textwidth]{Paper/figures/fig13Informe.png}
  \caption{Ingreso laboral por hora según posición ocupacional}
  \label{fig:14Informe}
\end{figure}

\textbf{Las posiciones de empleado de gobierno y patrón o empleador 
concentran los ingresos por hora más elevados del mercado laboral, 
mientras que el servicio doméstico y el trabajo por cuenta propia 
permanecen en los estratos más bajos de la distribución salarial.} 
En pesos constantes de 2025, el empleado de gobierno percibe \$25.785 
por hora y el patrón o empleador \$22.291, niveles que más que duplican el ingreso del empleado particular (\$11.417) y triplican el del trabajador por cuenta propia (\$8.031). El servicio doméstico, con \$6.506 por hora, es la categoría de menor remuneración. Comparado con 2010, todas las posiciones registran mejoras absolutas en términos reales, siendo los mayores avances absolutos los del empleado de gobierno (de \$19.738 a \$25.785) y el patrón o empleador (de \$16.212 a \$22.291), lo que amplía la distancia frente a las categorías de menor ingreso.

\begin{figure}[H]
  \centering
  \includegraphics[width=0.9\textwidth]{Paper/figures/fig14Informe.png}
  \caption{Crecimiento del ingreso por hora según posición ocupacional}
  \label{fig:15Informe}
\end{figure}

\textbf{A diferencia del patrón observado por nivel educativo, en la dimensión ocupacional son las posiciones de mayor ingreso las que registran el crecimiento relativo más alto, lo que indica una dinámica de divergencia salarial entre categorías ocupacionales.} Con base 2010 igual a 100, el patrón o empleador alcanza un índice de 137,5 en 2025, equivalente a un crecimiento anualizado del 2,15\%, el más alto entre todas las posiciones; le sigue el empleado particular con 133,2 y un 1,93\% anual, y el empleado de gobierno con 130,6 y un 1,8\% anual. Las categorías de menor remuneración crecen a ritmos inferiores: cuenta propia registra un índice de 129,3 (1,73\% anual) y servicio doméstico de 126,6 (1,59\% anual). Aunque las diferencias entre categorías no son drásticas en términos de tasa de crecimiento, el hecho de que las posiciones de mayor ingreso crezcan también más rápido implica que la brecha salarial absoluta entre ocupaciones se amplía, profundizando la segmentación del mercado laboral colombiano.

\subsection{Análisis por departamento}

\begin{figure}[H]
  \centering
  \includegraphics[width=0.9\textwidth]{Paper/figures/fig18Informe.png}
  \caption{Composición del empleo por departamento}
  \label{fig:16Informe}
\end{figure}

\begin{figure}[H]
  \centering
  \includegraphics[width=0.9\textwidth]{Paper/figures/fig21Informe.png}
  \caption{Composición del empleo por departamento}
  \label{fig:17Informe}
\end{figure}

\textbf{La distribución territorial del empleo en Colombia muestra una concentración persistente en Bogotá D.C. y Antioquia, aunque entre 2010 y 2025 se observa una leve redistribución hacia otras regiones del país.} Bogotá D.C. continúa siendo el principal polo de empleo nacional, aunque su participación se redujo 1,3 puntos porcentuales, pasando del 19,4\% al 18,1\% del total de ocupados; Antioquia, en cambio, ganó 1,3 puntos y se consolida como el segundo mercado laboral del país con el 14,4\%. Valle del Cauca y Atlántico mantienen su posición con el 8,8\% y el 5,5\% respectivamente, mientras que Cundinamarca emerge como uno de los territorios con mayor dinamismo, ganando 2,2 puntos porcentuales entre 2013 y 2025, lo que refleja el efecto de la expansión metropolitana de Bogotá hacia los municipios circundantes. Los departamentos de menor tamaño, agrupados en la categoría residual, aumentan ligeramente su participación conjunta en 1,3 puntos, lo que sugiere una incipiente pero aún débil desconcentración del empleo.

\begin{figure}[H]
  \centering
  \includegraphics[width=0.9\textwidth]{Paper/figures/fig16Informe.png}
  \caption{Ingreso laboral por hora según departamento: niveles}
  \label{fig:18Informe}
\end{figure}

\textbf{La jerarquía salarial entre departamentos es pronunciada y estable: Bogotá D.C. lidera consistentemente los ingresos por hora tanto en 2010 como en 2025, mientras que Córdoba, Nariño y Bolívar permanecen en los niveles más bajos de la distribución.} En 2025, el ingreso por hora en Bogotá D.C. asciende a \$15.904, seguido por Antioquia con \$11.791 y Cundinamarca con \$11.089; en el extremo opuesto, Córdoba registra \$6.197 y La Guajira \$5.520, menos de la mitad del ingreso bogotano. Comparado con 2010, todos los departamentos analizados muestran mejoras absolutas en términos reales, aunque las posiciones relativas en el ranking se mantienen en gran medida estables, con excepción de Risaralda y Caldas, que ascienden notoriamente. Esta persistencia de la jerarquía salarial territorial indica que las condiciones estructurales que determinan los ingresos en cada departamento ---composición sectorial, nivel educativo de la fuerza laboral, densidad económica--- cambian lentamente y no se modifican de manera sustancial en el horizonte de quince años analizado.

\begin{figure}[H]
  \centering
  \includegraphics[width=0.9\textwidth]{Paper/figures/fig17Informe.png}
  \caption{Ingreso laboral por hora según departamento: crecimiento 
2010--2025}
  \label{fig:19Informe}
\end{figure}

\textbf{Entre los 13 departamentos con información disponible desde 2010, Risaralda y Caldas lideran el crecimiento del ingreso real por hora, mientras que Santander registra el avance más modesto del grupo.} Con base 2010 igual a 100, Risaralda alcanza un índice de 152,2 en 2025, equivalente a un crecimiento anualizado del 2,84\%, el más alto del grupo; Caldas le sigue con 140,7 y un 2,3\% anual, y Bogotá D.C. con 136,7 y un 2,1\% anual. En contraste, Santander registra el crecimiento más bajo con un índice de 114,8 y apenas un 0,92\% anual, resultado llamativo dado que en términos de nivel absoluto se ubica entre los departamentos mejor remunerados. El resto del grupo ---Antioquia, Tolima, Nariño y Valle del Cauca--- se sitúa en una franja de crecimiento alto, entre el 1,68\% y el 1,95\% anual, coherente con el promedio nacional.

\begin{figure}[H]
  \centering
  \includegraphics[width=0.9\textwidth]{Paper/figures/fig20Informe.png}
  \caption{Ingreso laboral por hora según departamento: crecimiento 
2013--2025}
  \label{fig:20Informe}
\end{figure}

\textbf{La incorporación de 24 departamentos a partir de 2013 amplía el panorama regional y revela una heterogeneidad salarial aún más marcada, con Quindío y Risaralda liderando el crecimiento y La Guajira y Córdoba registrando caídas reales en el ingreso por hora.} En el ranking de niveles para 2025, Bogotá D.C. mantiene el primer lugar con \$15.904 por hora, y La Guajira cierra la distribución con \$5.520, una brecha de casi tres veces entre el departamento de mayor y menor ingreso. En términos de crecimiento entre 2013 y 2025, Quindío lidera con un índice de 146,3 y un 3,22\% anual, seguido de Risaralda (141,1; 2,91\%) y Caldas (131,2; 2,29\%), conformando un bloque del Eje Cafetero con desempeño salarial sobresaliente. En el extremo opuesto, La Guajira y Córdoba son los únicos departamentos con decrecimiento real del ingreso por hora entre 2013 y 2025, con tasas de $-$1,02\% y $-$0,71\% anual respectivamente, lo que profundiza su rezago relativo frente al resto del país y evidencia la existencia de dinámicas territoriales divergentes que la política laboral y regional debería atender de manera prioritaria.

\subsection{Análisis por sector económico}

\begin{figure}[H]
  \centering
  \includegraphics[width=0.9\textwidth]{Paper/figures/fig24Informe.png}
  \caption{Composición del empleo por sector económico}
  \label{fig:21Informe}
\end{figure}

\textbf{La estructura sectorial del empleo colombiano se ha transformado entre 2010 y 2025 con una contracción de los sectores tradicionales de mayor absorción laboral ---comercio y agricultura--- y una expansión de sectores de servicios de mayor valor agregado.} Comercio, el sector con mayor participación en ambos años, redujo su peso en 3 puntos porcentuales; agricultura lo hizo en 3,9 puntos, la mayor contracción relativa del período, y manufactura retrocedió 1,9 puntos. En sentido contrario, el sector de inmobiliarias y actividades profesionales registró el mayor avance con 3 puntos porcentuales adicionales, seguido de alojamiento y comida con 2,4 puntos y construcción y transporte con 1,2 puntos cada uno. Este patrón de terciarización del empleo ---con ganancia de sectores de servicios frente a la pérdida de primario y secundario--- es coherente con tendencias observadas en economías de ingreso medio en proceso de modernización productiva, aunque la contracción del empleo agrícola plantea interrogantes sobre la absorción de trabajadores rurales desplazados hacia sectores urbanos de baja productividad.

\begin{figure}[H]
  \centering
  \includegraphics[width=0.9\textwidth]{Paper/figures/fig22Informe.png}
  \caption{Ingreso laboral por hora según sector económico: niveles}
  \label{fig:22Informe}
\end{figure}

\textbf{La heterogeneidad salarial entre sectores es extrema: en 2025, los trabajadores de información y comunicaciones, administración pública y financieras perciben ingresos por hora que más que triplican los de agricultura, hogares empleadores y alojamiento y comida.} En pesos constantes de 2025, los tres sectores de mayor remuneración son información y comunicaciones (\$21.868), administración pública (\$21.557) y financieras (\$21.538), mientras que en el extremo inferior se ubican agricultura (\$6.080), hogares empleadores (\$6.509) y alojamiento y comida (\$6.962). Comparado con 2010, todos los sectores registran mejoras en términos reales, siendo el avance más notable el de información y comunicaciones, que pasó de \$8.735 a \$21.868. La brecha absoluta entre el sector mejor y peor remunerado se amplió sustancialmente en el período, pasando de aproximadamente \$10.000 en 2010 a más de \$15.000 en 2025, lo que evidencia una creciente segmentación salarial de origen sectorial.

\begin{figure}[H]
  \centering
  \includegraphics[width=0.9\textwidth]{Paper/figures/fig23Informe.png}
  \caption{Crecimiento del ingreso por hora según sector económico}
  \label{fig:23Informe}
\end{figure}

\textbf{El sector de información y comunicaciones registra un crecimiento del ingreso real por hora sin precedentes entre 2010 y 2025, con un índice de 250,3 que lo distancia ampliamente del resto de sectores y refleja la digitalización acelerada de la economía colombiana.} Con base 2010 igual a 100, información y comunicaciones alcanza 250,3 en 2025, equivalente a un crecimiento anualizado del 6,31\%, más del doble que el segundo sector de mayor crecimiento, administración pública, con 138 y un 2,17\% anual. Le siguen manufactura (128,8; 1,7\% anual), agricultura (127; 1,61\% anual) y hogares empleadores (126,8; 1,6\% anual). En el extremo opuesto, minas (110,1; 0,65\% anual) y servicios públicos (105,5; 0,36\% anual) muestran el crecimiento más modesto, posiblemente asociado a la volatilidad de los precios internacionales de commodities y a la rigidez tarifaria regulada, respectivamente. El patrón general indica que los sectores intensivos en conocimiento y tecnología han concentrado las mayores ganancias salariales del período.

\begin{figure}[H]
  \centering
  \includegraphics[width=0.9\textwidth]{Paper/figures/fig27Informe.png}
  \caption{Composición interna de las subramas seleccionadas}
  \label{fig:24Informe}
\end{figure}

\textbf{Al interior de los sectores, la composición por subramas revela transformaciones significativas: el subsector de software, informática y sistemas más que duplicó su participación dentro del sector de información y comunicaciones, consolidándose como el núcleo del dinamismo tecnológico del país.} Entre las subramas de mayor ingreso en 2025, software, informática y sistemas pasó de representar el 17,7\% del sector de información y comunicaciones en 2010 al 46,5\% en 2025, mientras que servicios de información y datos creció del 0,4\% al 4,8\%. En el sector financiero, seguros y auxiliares financieros ganó 10 puntos porcentuales dentro de su rama, pasando del 26,7\% al 36,7\%. Entre las subramas de menor ingreso, agricultura y ganadería concentra prácticamente la totalidad del empleo agrícola (95,6\% en 2025), y hogares como empleadores mantiene el 100\% de su sector por definición. Estos cambios composicionales internos refuerzan la lectura de que la modernización del mercado laboral colombiano se está canalizando fundamentalmente a través de la expansión de actividades digitales y de conocimiento dentro de sectores ya dinámicos.

\begin{figure}[H]
  \centering
  \includegraphics[width=0.9\textwidth]{Paper/figures/fig25Informe.png}
  \caption{Ingreso laboral por hora según subrama detallada: niveles 
y crecimiento}
  \label{fig:25Informe}
\end{figure}

\begin{figure}[H]
  \centering
  \includegraphics[width=0.9\textwidth]{Paper/figures/fig26Informe.png}
  \caption{Ingreso laboral por hora según subrama detallada: niveles 
y crecimiento}
  \label{fig:26Informe}
\end{figure}

\textbf{Las cinco subramas de mayor ingreso en 2025 pertenecen todas al universo digital, científico y financiero, con servicios de información y datos liderando tanto en nivel (\$32.275 por hora) como en crecimiento (7,46\% anual), en tanto que cuero y calzado es la única subrama analizada con caída real del ingreso por hora en el período.} En términos de nivel, las subramas mejor remuneradas en 2025 son servicios de información y datos (\$32.275), software, informática y sistemas (\$28.061), investigación y desarrollo (\$27.478), profesionales, científicas y técnicas (\$23.057) y seguros y auxiliares financieros (\$22.254); todas ellas con crecimientos reales sustanciales respecto a 2010. En el extremo opuesto, las subramas de menor ingreso ---pesca y acuicultura (\$4.202), agricultura y ganadería (\$4.841), cuero y calzado (\$5.228), silvicultura (\$5.403) y hogares como empleadores (\$5.132)--- permanecen en niveles que representan entre una séptima y una sexta parte del ingreso de las subramas líderes. En materia de crecimiento, servicios de información y datos alcanza un índice de 294,3 (7,46\% anual), software e informática de 169,5 (3,58\% anual) e investigación y desarrollo de 168,2 (3,53\% anual), mientras que cuero y calzado registra el único decrecimiento real con un índice de 91 ($-$0,62\% anual), lo que subraya la presión competitiva que enfrenta este subsector manufacturero de baja tecnología.

\subsection{Análisis por tamaño de empresa}

\begin{figure}[H]
  \centering
  \includegraphics[width=0.9\textwidth]{Paper/figures/fig30Informe.png}
  \caption{Composición del empleo según tamaño de empresa}
  \label{fig:27Informe}
\end{figure}

\textbf{El mercado laboral colombiano está dominado por el trabajo unipersonal y las microempresas, aunque entre 2010 y 2025 se observa un desplazamiento neto hacia las empresas de mayor tamaño, particularmente las de 101 o más trabajadores.} En 2025, el 34,8\% de los ocupados trabaja solo ---es decir, sin empleados a cargo---, y el 13,8\% en empresas de 2 a 3 personas, de modo que casi la mitad de la fuerza laboral se concentra en unidades productivas de menos de cuatro personas. Sin embargo, respecto a 2010, estas dos categorías han perdido participación: los trabajadores solos cayeron 2,7 puntos porcentuales y los de empresas de 2 a 3 personas, 5,1 puntos. En contraste, las empresas de 101 o más trabajadores ganaron 5,1 puntos porcentuales, pasando del 20,0\% al 25,1\%, el avance más significativo del período. Los tramos intermedios ---de 6 a 100 trabajadores--- registran ganancias modestas de entre 0,2 y 1 punto porcentual cada uno. Este proceso de relativa concentración del empleo en empresas grandes es consistente con la tendencia hacia la formalización documentada en secciones anteriores, dado que las empresas de mayor tamaño tienden a ofrecer condiciones laborales más protegidas.

\begin{figure}[H]
  \centering
  \includegraphics[width=0.9\textwidth]{Paper/figures/fig28Informe.png}
  \caption{Composición del empleo según tamaño de empresa}
  \label{fig:28Informe}
\end{figure}

\textbf{Existe una relación positiva y clara entre el tamaño de la empresa y el ingreso por hora de sus trabajadores, relación que se mantiene estable entre 2010 y 2025 y que implica que los ocupados en empresas grandes ganan más del doble que quienes trabajan solos.} En pesos constantes de 2025, los trabajadores en empresas de 101 o más personas perciben \$17.076 por hora, frente a \$6.674 de quienes trabajan solos y \$7.339 de quienes lo hacen en unidades de 2 a 3 personas. El gradiente salarial es prácticamente monotónico a lo largo de todos los tramos de tamaño: a mayor número de trabajadores en la empresa, mayor ingreso por hora en todos los años observados. Comparado con 2010, todos los segmentos registran mejoras absolutas en términos reales, siendo el avance más pronunciado en las empresas de 101 o más trabajadores, que pasaron de \$13.963 a \$17.076, y en las de 51 a 100, que pasaron de \$10.609 a \$13.448. Este patrón refuerza la idea de que la inserción en empresas de mayor tamaño constituye uno de los determinantes más robustos del nivel salarial en el mercado laboral colombiano.

\begin{figure}[H]
  \centering
  \includegraphics[width=0.9\textwidth]{Paper/figures/fig29Informe.png}
  \caption{Crecimiento del ingreso por hora según tamaño de empresa}
  \label{fig:29Informe}
\end{figure}

\textbf{A diferencia de lo que ocurre con los niveles, el crecimiento relativo del ingreso por hora no sigue un patrón lineal respecto al tamaño de la empresa: son las empresas de tamaño mediano-bajo (11 a 19 trabajadores) y mediano-alto (51 a 100) las que registran los mayores ritmos de crecimiento real entre 2010 y 2025.} Con base 2010 igual a 100, las empresas de 11 a 19 trabajadores alcanzan un índice de 131 en 2025, equivalente al mayor crecimiento anualizado del grupo con un 1,82\%; le siguen las de 51 a 100 trabajadores con 126,8 y un 1,59\% anual. En contraste, las empresas de 4 a 5 trabajadores registran el crecimiento más bajo con un índice de 115,8 y un 0,98\% anual, y las de 20 a 30 con 116,3 y un 1,01\% anual. Las empresas más grandes (101 o más), pese a concentrar los mayores ingresos absolutos, crecen a un ritmo de apenas el 1,35\% anual, por debajo del promedio nacional. Esta dinámica sugiere que ciertos segmentos del tejido empresarial mediano están mejorando sus condiciones salariales de manera más acelerada, aunque sin alterar la jerarquía estructural en la que las empresas grandes siguen liderando los niveles de remuneración.

\newpage
\section{Conclusión}

\textbf{El mercado laboral colombiano de 2025 es simultáneamente más grande, más formal y mejor remunerado que el de 2010, pero estas mejoras agregadas encubren una heterogeneidad profunda que limita su alcance en términos de equidad y bienestar generalizado.} A lo largo de las secciones anteriores se ha documentado que el crecimiento del empleo y del ingreso real por hora ha sido sostenido en el período analizado, con una recuperación postpandemia notable que en apenas cuatro años no solo repuso los empleos perdidos sino que los superó ampliamente. Sin embargo, esta narrativa positiva debe leerse con cautela, pues los promedios nacionales agregan realidades muy distintas según la dimensión de análisis que se adopte.

\textbf{La formalidad laboral es la dimensión que más condiciona las oportunidades de ingreso en Colombia, y pese a los avances registrados, más de la mitad de los trabajadores permanece desprotegida.} La brecha salarial entre formales e informales ---de 2,43 veces en 2025--- no solo no se cerró sino que se amplió respecto a 2010, lo que significa que la formalización, aunque real en términos de flujos, no ha redistribuido proporcionalmente los beneficios del crecimiento hacia quienes permanecen en la informalidad. Acelerar la transición hacia el empleo formal es, por tanto, la palanca de política con mayor potencial redistributivo dentro del mercado laboral.

\textbf{La digitalización de la economía es el fenómeno sectorial más disruptivo del período y su impacto sobre los ingresos laborales es de una magnitud sin precedentes en la historia reciente del mercado de trabajo colombiano.} El sector de información y comunicaciones triplicó su ingreso real por hora entre 2010 y 2025, y subramas como servicios de información y datos casi triplicaron su nivel salarial con un crecimiento anualizado del 7,46\%. Este dinamismo, sin embargo, beneficia a un segmento aún reducido de la fuerza laboral, altamente cualificado y concentrado en los principales centros urbanos, lo que hace de la expansión del acceso a la educación superior y a las competencias digitales una prioridad de política para ampliar los beneficios de la transformación tecnológica.

\textbf{Las brechas territoriales constituyen uno de los rasgos más persistentes y preocupantes del mercado laboral colombiano, con departamentos como La Guajira y Córdoba que no solo parten de niveles salariales bajos sino que registraron caídas reales en el ingreso por hora entre 2013 y 2025.} La brecha entre Bogotá D.C. y La Guajira ---de casi tres veces en el ingreso por hora--- refleja diferencias estructurales en composición sectorial, capital humano, densidad económica e infraestructura que no se corrigen espontáneamente con el crecimiento agregado y que requieren intervenciones de política regional activa y sostenida.

\textbf{La convergencia educativa de la fuerza laboral es una de las transformaciones más alentadoras del período, aunque su traducción en reducción de la desigualdad salarial absoluta es aún marginal.} El hecho de que un tercio de los ocupados cuente hoy con educación superior, frente a menos de un cuarto en 2010, representa un cambio estructural de largo alcance. No obstante, dado que los retornos educativos en términos de ingreso por hora siguen siendo muy elevados y que la brecha absoluta entre niveles no se ha reducido, la expansión educativa por sí sola no garantiza una distribución más equitativa de los ingresos laborales si no va acompañada de políticas que amplíen la demanda de trabajo cualificado en sectores de alta productividad.

\textbf{En síntesis, el período 2010--2025 configura un mercado laboral colombiano en transición: más moderno en su composición, más conectado a la economía digital y con una base formal más amplia, pero aún atravesado por brechas estructurales de género, formalidad, territorio y cualificación que definen trayectorias de ingreso radicalmente distintas según el perfil del trabajador.} Los avances documentados en este informe son reales y significativos, pero su consolidación y profundización requieren políticas laborales, educativas y regionales que actúen simultáneamente sobre los múltiples ejes de desigualdad identificados, con especial atención a los grupos y territorios que han quedado rezagados en la distribución de los beneficios del crecimiento.


\bibliographystyle{apalike}
\bibliography{Paper/references}

\end{document}
